% Part: first-order-logic
% Chapter: tableaux
% Section: proof-theoretic-notions

\documentclass[../../../include/open-logic-section]{subfiles}

\begin{document}

\iftag{FOL}
      {\olfileid[id]{fol}{tab}{ptn}}
      {\olfileid[id]{pl}{tab}{ptn}}

\olsection{Gagasan Teori-Bukti}

\begin{editorial}
Bagian ini menghimpun definisi relasi keterderivasian dan konsistensi
untuk tableau.
\end{editorial}

\begin{explain}
Sebagaimana kita telah mendefinisikan sejumlah gagasan semantik penting
(validitas, perikutan, keterpenuhan), kini kita mendefinisikan
\emph{gagasan teori-bukti} yang bersesuaian. Gagasan-gagasan ini tidak
didefinisikan dengan mengacu pada terpenuhinya !!{sentence} di dalam
\iftag{FOL}{!!{structure}}{!!{valuation}}, melainkan dengan mengacu pada keberadaan !!{tableau}s
tertutup tertentu. Penemuan penting menunjukkan bahwa gagasan-gagasan
ini berimpit. Keberimpitan itu adalah isi \emph{teorema kesahihan} dan
\emph{teorema kelengkapan}.
\end{explain}


\begin{defn}[Teorema]
!!^a{sentence}~$!A$ adalah sebuah \emph{teorema} jika terdapat
!!{tableau} tertutup untuk~$\sFmla{\False}{!A}$. Kita menulis
$\Proves !A$ jika $!A$ merupakan teorema dan $\Proves/ !A$ jika bukan.
\end{defn}

\begin{defn}[!!^{derivability}]
!!^a{sentence} $!A$ \emph{!!{derivable} dari} suatu himpunan
!!{sentence}~$\Gamma$, $\Gamma \Proves !A$, jika dan hanya jika terdapat
himpunan berhingga $\{!B_1, \dots, !B_n\} \subseteq \Gamma$ dan
!!{tableau} tertutup untuk himpunan
\[
\{
  \sFmla{\False}{!A},
  \sFmla{\True}{!B_1}, \dots,
  \sFmla{\True}{!B_n}
  \}.
\]
Jika $!A$ tidak !!{derivable} dari $\Gamma$, kita menulis
$\Gamma \Proves/ !A$.
\end{defn}

\begin{defn}[Konsistensi]
Suatu himpunan !!{sentence}~$\Gamma$ \emph{tak konsisten} jika dan hanya
jika terdapat himpunan berhingga $\{!B_1, \dots, !B_n\} \subseteq
\Gamma$ dan !!{tableau} tertutup untuk himpunan
\[
\{
  \sFmla{\True}{!B_1}, \dots,
  \sFmla{\True}{!B_n}
  \}.
\]
Jika $\Gamma$ tidak tak konsisten, kita menyebutnya \emph{konsisten}.
\end{defn}

\begin{prop}[Refleksivitas]
\ollabel{prop:reflexivity}
Jika $!A \in \Gamma$, maka $\Gamma \Proves !A$.
\end{prop}

\begin{proof}
Jika $!A \in \Gamma$, maka $\{!A\}$ merupakan himpunan bagian berhingga
dari~$\Gamma$, dan !!{tableau}
\begin{oltableau}
  [\sFmla{\False}{\formula{A}}, just = \TAss
    [\sFmla{\True}{\formula{A}}, just = \TAss,close]
  ]
\end{oltableau}
tertutup.
\end{proof}


\begin{prop}[Monotonisitas]
\ollabel{prop:monotonicity}
Jika $\Gamma \subseteq \Delta$ dan $\Gamma \Proves !A$, maka $\Delta
\Proves !A$.
\end{prop}

\begin{proof}
Setiap himpunan bagian berhingga dari~$\Gamma$ juga merupakan himpunan
bagian berhingga dari~$\Delta$.
\end{proof}

\begin{prop}[Transitivitas]
\ollabel{prop:transitivity}
Jika $\Gamma \Proves !A$ dan $\{!A\} \cup
\Delta \Proves !B$, maka $\Gamma \cup \Delta \Proves !B$.
\end{prop}

\begin{proof}
Jika $\{!A\} \cup \Delta \Proves !B$, maka terdapat himpunan bagian
berhingga $\Delta_0 = \{!C_1, \dots, !C_n\} \subseteq \Delta$
sedemikian sehingga
\begin{align*}
\{\sFmla{\False}{!B}, & \sFmla{\True}{!A}, \sFmla{\True}{!C_1},
\dots, \sFmla{\True}{!C_n}\}
\intertext{memiliki !!{tableau} tertutup. Jika $\Gamma \Proves !A$,
  maka terdapat $\{!D_1, \dots, !D_m\} \subseteq \Gamma$ sedemikian
  sehingga}
\{\sFmla{\False}{!A}, & \sFmla{\True}{!D_1},
\dots, \sFmla{\True}{!D_m}\}
\end{align*}
memiliki !!{tableau} tertutup.

Sekarang perhatikan !!{tableau} dengan asumsi-asumsi
\[
\sFmla{\False}{!B},
\sFmla{\True}{!C_1}, \dots, \sFmla{\True}{!C_n},
\sFmla{\True}{!D_1}, \dots, \sFmla{\True}{!D_m}.
\]
Terapkan aturan \Cut{} pada~$!A$. Penerapan ini menghasilkan dua cabang:
yang satu memuat $\sFmla{\True}{!A}$, sedangkan yang lain memuat
$\sFmla{\False}{!A}$. Jadi, pada cabang pertama, semua anggota
\[
\{\sFmla{\False}{!B}, \sFmla{\True}{!A},
\sFmla{\True}{!C_1}, \dots, \sFmla{\True}{!C_n}\}
\]
tersedia. Karena terdapat !!{tableau} tertutup untuk asumsi-asumsi ini,
kita dapat menempelkannya pada cabang tersebut; setiap cabang yang
melewati $\sFmla{\True}{!A}$ tertutup. Pada cabang lainnya, semua anggota
\[
\{\sFmla{\False}{!A}, \sFmla{\True}{!D_1}, \dots,
\sFmla{\True}{!D_m}\}
\]
tersedia, sehingga kita juga dapat melengkapi sisi lainnya untuk
memperoleh !!{tableau} tertutup. Hal ini menunjukkan
$\Gamma \cup \Delta \Proves !B$.
\end{proof}

Perhatikan bahwa, khususnya, ini berarti bahwa jika $\Gamma \Proves !A$
dan $!A \Proves !B$, maka $\Gamma \Proves !B$. Diperoleh pula bahwa jika
$!A_1, \dots, !A_n \Proves !B$ dan $\Gamma \Proves !A_i$ untuk
setiap~$i$, maka $\Gamma \Proves !B$.

\begin{prop}
\ollabel{prop:incons}
$\Gamma$ tak konsisten jika dan hanya jika $\Gamma \Proves !A$ untuk
setiap !!{sentence}~$!A$.
\end{prop}

\begin{proof}
Latihan.
\end{proof}

\tagprob{FOL}
\begin{prob}
Buktikan \olref[fol][tab][ptn]{prop:incons}.
\end{prob}
\tagendprob

\tagprob{notFOL}
\begin{prob}
Buktikan \olref[pl][tab][ptn]{prop:incons}.
\end{prob}
\tagendprob

\begin{prop}[Kekompakan]
  \ollabel{prop:proves-compact}
  \begin{enumerate}
  \item Jika $\Gamma \Proves !A$, maka terdapat himpunan bagian berhingga
    $\Gamma_0 \subseteq \Gamma$ sedemikian sehingga
    $\Gamma_0 \Proves !A$.
  \item Jika setiap himpunan bagian berhingga dari~$\Gamma$ konsisten,
    maka $\Gamma$ konsisten.
  \end{enumerate}
\end{prop}

\begin{proof}
  \begin{enumerate}
    \item Jika $\Gamma \Proves !A$, maka terdapat himpunan bagian
      berhingga $\Gamma_0 = \{!B_1, \dots, !B_n\} \subseteq \Gamma$
      dan !!{tableau}
      tertutup untuk
      \[
      \{\sFmla{\False}{!A}, \sFmla{\True}{!B_1}, \dots,
      \sFmla{\True}{!B_n}\}
      \]
      !!^{tableau} ini juga menunjukkan $\Gamma_0 \Proves !A$.
    \item Jika $\Gamma$ tak konsisten, maka untuk suatu himpunan bagian
      berhingga $\Gamma_0 = \{!B_1, \dots, !B_n\}$ terdapat
      !!{tableau} tertutup untuk
      \[
      \{\sFmla{\True}{!B_1}, \dots, \sFmla{\True}{!B_n}\}
      \]
      !!^{tableau} tertutup ini menunjukkan bahwa $\Gamma_0$ tak
      konsisten.
  \end{enumerate}
\end{proof}

\end{document}
